\section{Introducción: llega la era de los Agent}

Este episodio, presentado por Greg Isenberg, tiene como invitado a Boris, de Anthropic —creador de Claude Code e impulsor central de Claude Co-Work—, en una entrevista. Boris no solo muestra el flujo de operación de Co-Work, sino que también comparte su propio flujo de trabajo y mejores prácticas al usar Claude Code. Este apunte organiza sistemáticamente el contenido de la conversación en cuatro ejes: posicionamiento del producto, forma de uso, mecanismos de seguridad y flujos de trabajo eficientes.

\begin{importantbox}{Idea central}
La esencia de Claude Co-Work es empaquetar la capacidad de Agent de Claude Code en una interfaz que cualquier persona pueda usar; por debajo corre exactamente el mismo Claude Agent SDK que Claude Code. Co-Work puede operar archivos, controlar el navegador y conectarse a cualquier herramienta mediante MCP: no es solo un chat, sino un Agent en el sentido pleno del término.
\end{importantbox}

\begin{knowledgebox}{¿Qué es agentic AI?}
Boris señala que la palabra ``agentic'' se ha usado de forma abusiva en la industria. Su significado central es: \textbf{la IA puede tomar acciones (take action)}, es decir, no se limita a generar texto o buscar en la web, sino que es capaz de usar las herramientas de la computadora e interactuar con el mundo real. Desde muy temprano, Anthropic definió coding, tool use y computer use como los tres ejes centrales de entrenamiento del modelo.
\end{knowledgebox}

\subsection{Resumen de la sección}

Co-Work es el producto de Agent que Anthropic lanzó para usuarios sin conocimientos técnicos: presenta toda la capacidad de Agent de Claude Code a través de una interfaz de escritorio sencilla. Comprender el concepto central de que ``Agent significa poder tomar acciones'' es la base para entender todas las funciones de Co-Work.

